\section{Related Works}
\label{sec:related}

For the elaboration of this study, comparative works on time series databases were reviewed,
covering both functionality analyses and performance benchmarks. The search was conducted
using keywords such as ``data streaming'', ``timeseries'', and ``timeseries databases'',
with results ordered by citation count to prioritize more established and developed analyses.
From this search, 23 results were selected for initial reading to verify their relevance to
the objectives of this study.

Among the 23 works, some did not present sufficient direct relation for an in-depth analysis,
though they provided relevant context. \cite{zubaroglu2021data}, for example, presents a list
of datasets applicable to time series database analyses. \cite{Gaber2005Mining} contributed
relevant information about the algorithms underlying TSDB implementations, despite not being
directly focused on comparative evaluation.

From the initial reading, six works were selected for detailed study. This chapter is organized
in three sections. The \textit{Functionality Analysis} section covers works that evaluate
technical and operational characteristics of TSDB systems. The \textit{Performance Analysis}
section covers comparative benchmark studies. The \textit{Comparative Analysis} section
synthesizes and critiques all selected works through a comparative table, highlighting patterns,
gaps, and trends identified in the literature.

\subsection{Functionality Analysis}

\cite{Petre2019} present a maturity-model analysis of six open-source time series databases:
InfluxDB, Graphite, RRDTool, Prometheus, OpenTSDB, and TimescaleDB. The work evaluates 18
quantitative and qualitative attributes divided into functional requirements (apparent to the
end user, such as specific capabilities) and non-functional requirements (not directly
apparent but impactful, such as compression strategies and query optimization mechanisms).
Evaluated attributes include active development, development velocity, market maturity, data
replication, and library support. The objective is to assist in the selection of the most
adequate TSDB for business needs, considering both technical and operational aspects. The
result identified InfluxDB as the leading solution in the analysis.

\subsection{Performance Analysis}

Five works in the format of controlled benchmarks or real-scenario comparisons were selected
for detailed analysis.

\cite{lima2024avaliando} employed insertion tests (ranging from 1{,}000 to 5 million records)
and aggregate queries with synthetic data (timestamp, location, temperature, and humidity
measurements) to compare DBMS performance using Apache JMeter. The same data structure was
replicated across all databases. Tests were carried out in a Docker-based environment on a
machine with AMD Ryzen 5 4500 (6 cores), 16\,GB RAM, and SSD storage, testing at volumes of
500{,}000 and 5 million records. Reproducibility was partial, as Docker containers were used
but with incomplete details for full replication. InfluxDB was the standout performer in
the study.

\cite{shah2022performance} compared the performance of InfluxDB~v2.1, TimescaleDB~v2.6,
Druid~v0.22.1, and Cassandra~v3.11.12 in insertion operations and eight query types (point
queries, aggregations, and others), evaluating load time, space consumption, and latency.
Both real datasets (IoT data, financial data, NYC Taxi) and synthetic datasets (generated via
a Python time-series generator) were used. The test environment consisted of an Intel Core
i5-9500 (6 cores), 8\,GB RAM, 500\,GB SSD, and Ubuntu~20.04. InfluxDB stood out not only
in query performance but also in space consumption, ranking second only to Druid. The study
was classified as non-reproducible, as no scripts or configurations were shared.

\cite{naqvi2017influxdb} compared InfluxDB, SQL Server, Cassandra, Elasticsearch, and OpenTSDB
on three primary metrics: data ingestion rate, storage requirements, and average query response
time. Both synthetic data (simulating server metrics) and a real dataset of NYC taxi trips
(2009--2017) were used, processed with native DBMS tools and custom Java scripts. Tests were
performed on an Intel Core i5-6300HQ (4 cores), 16\,GB RAM, and SSD, using specific versions
of each database (InfluxDB~1.3, Elasticsearch~5.6.3). The study concluded that InfluxDB is
ideal for high-velocity temporal data ingestion scenarios such as IoT and monitoring, but may
not be adequate for applications requiring complex JOIN operations or frequent updates.
Reproducibility was classified as non-existent, with no scripts or configurations shared.

\cite{pereira2023mulletbench} evaluated the performance of InfluxDB~2.7 and Apache IoTDB~1.1.1
in three scenarios: cloud-only, 2~edge~+~cloud, and 4~edge~+~cloud, testing insertion,
query, and mixed workloads, each executed three times after a system reboot. The tools used
included MulletBench for the benchmark, Ansible and Docker for orchestration, PCP for
monitoring, and tcconfig for network simulation. Tests were conducted on server nodes with
Intel~i5-9500 (16\,GB RAM, NVMe SSD) and client nodes with Intel~i3-4170 (16\,GB RAM,
SATA SSD) over a Gigabit switched network. This is the only work in the reviewed literature
classified as fully reproducible, with all configurations and code publicly available on GitHub.
Apache IoTDB was the standout performer, outperforming InfluxDB across all tested scenarios.

\cite{wang2023iotdb} compared IoTDB~v0.13.0 against InfluxDB, TimescaleDB, KairosDB, Parquet,
and ORC in write and read operations, storage cost, and query types including range queries,
aggregations, and downsampling, measuring throughput, latency, and response time.
Both real industrial datasets (vehicle, train, and machine data) and synthetic datasets
generated by TSBS were used. Tests were conducted in standalone (Intel i7-10700, 32\,GB RAM,
10\,TB HDD) and distributed (Alibaba Cloud with Xeon Platinum, 256\,GB RAM) environments,
with a replication factor of 3 and temporal partitioning in the distributed setup. While
configuration details are described in the paper, no complete reproduction guide was provided,
leading to a partial reproducibility classification. Apache IoTDB demonstrated superiority
in both local and distributed environments, especially at large industrial scales involving
trillions of data points.

\subsection{Comparative Analysis}

Table~\ref{tab:related} aggregates the key attributes of the five performance studies, enabling
a comparative analysis across dimensions most relevant to benchmark methodology.

\begin{table*}[t]
  \centering
  \caption{Comparison of related TSDB benchmark studies.}
  \label{tab:related}
  \begin{tabular}{lllllll}
    \hline
    \textbf{Author} & \textbf{Year} & \textbf{Databases} & \textbf{Reprod.} & \textbf{Datasets} & \textbf{Tools} & \textbf{Leader} \\
    \hline
    \cite{lima2024avaliando}       & 2024 & SQL + NoSQL & Partial & Synthetic          & Apache JMeter                 & InfluxDB \\
    \cite{pereira2023mulletbench}  & 2023 & NoSQL       & Full    & Synthetic          & MulletBench, Ansible, Docker  & IoTDB    \\
    \cite{wang2023iotdb}           & 2023 & SQL + NoSQL & Partial & Real + Synthetic   & TSBS                          & IoTDB    \\
    \cite{shah2022performance}     & 2022 & SQL + NoSQL & None    & Real + Synthetic   & timeseries-generator (Python) & InfluxDB \\
    \cite{naqvi2017influxdb}       & 2017 & SQL + NoSQL & None    & Real + Synthetic   & Telegraf, Kapacitor            & InfluxDB \\
    \hline
  \end{tabular}
\end{table*}

The publication years span from 2017 to 2024. The majority of studies are recent, with
particular density in 2023, reflecting growing interest in TSDB evaluation. The presence of
studies from 2017 highlights that this is not a newly emerging concern, yet the rapid evolution
of the databases themselves means that older findings may no longer apply to current versions.

\textbf{Database coverage.} InfluxDB is present in every reviewed study, serving as the de
facto reference point for TSDB comparisons, even in works where another system is the
primary subject, as in~\cite{wang2023iotdb}. The presence of TimescaleDB in three of five
studies reflects its growing relevance as a SQL-compatible alternative. The most notable
trend is the emergence of Apache IoTDB: appearing only in the two most recent studies
(2023), it leads both of them~\citep{pereira2023mulletbench,wang2023iotdb}, suggesting performance improvements in recent versions, selection bias in recent benchmarks
toward IoT workloads, or growing recognition of its strengths for this domain; the
available data do not allow distinguishing among these explanations.

\textbf{Reproducibility.} This is the most significant methodological gap in the reviewed
literature. Of five performance studies, only one, \cite{pereira2023mulletbench},
is fully reproducible, with all code and configurations publicly available on GitHub.
Two studies (\cite{lima2024avaliando} and \cite{wang2023iotdb}) offer partial reproducibility
through Docker containers or described configurations, while the remaining two
(\cite{shah2022performance} and \cite{naqvi2017influxdb}) provide no means of replication.
This limits independent validation of results and makes it difficult to attribute performance
differences to database architecture rather than experimental conditions.

\textbf{Datasets.} All five studies use synthetic datasets, reflecting the need for
controlled, volume-scalable data. Three studies additionally include real datasets such as
NYC Taxi data, IoT sensor readings, and industrial vehicle data~\citep{shah2022performance,
naqvi2017influxdb,wang2023iotdb}, providing practical relevance alongside experimental
control. The TSBS (Time Series Benchmark Suite) appears in the most recent studies,
suggesting a potential standardization trend in the community.

\textbf{Hardware configurations.} Tested environments vary considerably: processors range
from Intel~i3 to Intel~i7, RAM from 8\,GB to 32\,GB (plus cloud with 256\,GB in one
case~\citep{wang2023iotdb}), and storage from SATA SSD to NVMe and HDD. The majority of
tests used mid-range machines, consistent with a focus on realistic deployment scenarios.
The inclusion of cloud testing in more recent studies~\citep{pereira2023mulletbench,wang2023iotdb}
reflects growing interest in edge-to-cloud architectures.

\textbf{Outcome trends.} InfluxDB leads in three of five studies, all from 2017--2022.
Apache IoTDB leads in both 2023 studies. This temporal pattern is notable: it may indicate
performance improvements in IoTDB's recent versions, or it may reflect the fact that
more recent benchmarks include IoT-specific workloads where IoTDB's native data model
provides a structural advantage. The shift does not necessarily invalidate InfluxDB for
other workloads; it does, however, motivate the inclusion of both databases in any
contemporary comparative evaluation.

In summary, the reviewed literature reveals three key gaps that the present work addresses:
(1) insufficient reproducibility: only one of five studies is fully replicable; (2) limited
coverage of IoTDB alongside InfluxDB and TimescaleDB in the same controlled benchmark; and
(3) the absence of a systematic analysis of how methodology choices (resource isolation,
database tuning) affect reported conclusions. This study contributes to closing all three gaps.
